% Complete TeX conversion of THREE_LINE_IMAGE_PROOF.md, Sections 1--5.
% Retains the original F=tau from es_niemeier_wilson_cyclic.tex.
\section{The exact cubic image of the original common lattice}
\label{d4i:sec:complete-image}

Retain the coordinates $a_0,a_1,a_2,a_3$ and the original lattice
\[
 D_4=\{a\in\mathbb Z^4:a_0+a_1+a_2+a_3\in2\mathbb Z\}.
 \label{d4i:eq:original-domain}\tag{D4I1}
\]
The polynomial already used in \cref{wcl:sec:common-cubic} is
\[
 P(a)=a_0\bigl(a_0^2-3(a_1^2+a_2^2+a_3^2)\bigr).
 \label{d4i:eq:original-polynomial}\tag{D4I2}
\]
Here and throughout this section $F$ is the original triality cubic
$F=\tau$, with its ordered octonionic arguments as in
\cref{wcl:sec:construction,wcl:sec:common-cubic}. In particular, on the
literal common lattice of \cref{wcl:thm:literal-intersection},
\[
 \mathcal I=\{\sqrt2(a,a,a):a\in D_4\},\qquad
 F(\sqrt2(a,a,a))=2\sqrt2\,P(a).
 \label{d4i:eq:retained-cubic}\tag{D4I3}
\]
No coordinate permutation, scaling of $P$, or alteration of the lattice
is made. The witnesses below are actual points of this original lattice.
The image proof is elementary and requires no external arithmetic
existence theorem.

\subsection{Necessary restrictions in the original coordinates}
\label{d4i:sec:necessary}

For every integer $b$, one has $b^2\equiv b\pmod2$. Consequently,
for $a\in D_4$,
\[
 a_0^2+a_1^2+a_2^2+a_3^2
 \equiv a_0+a_1+a_2+a_3\equiv0\pmod2.
\]
Writing $s=a_1^2+a_2^2+a_3^2$, the factor $a_0^2-3s$ is congruent
to $a_0^2+s\equiv0\pmod2$. Thus $P(a)$ is even.

Modulo three the original polynomial satisfies $P(a)\equiv a_0^3$.
If $3\mid P(a)$, the field $\mathbb Z/3\mathbb Z$ has no nonzero
element with zero cube, so $3\mid a_0$. Substitution of $a_0=3b$
into the original formula gives
\[
 P(a)=3b(9b^2-3s)=9b(3b^2-s).
 \label{d4i:eq:three-divisibility}\tag{D4I4}
\]
Therefore
\[
 P(a)\in2\mathbb Z,\qquad
                 3\mid P(a)\ \Longrightarrow\ 9\mid P(a).
 \label{d4i:eq:necessary-restrictions}\tag{D4I5}
\]

\subsection{Three explicit affine lines and their exact values}
\label{d4i:sec:three-lines}

Define the common direction
\[
 v=(-3,1,1,1).
 \label{d4i:eq:common-direction}\tag{D4I6}
\]
Its coordinate sum is zero, and direct substitution gives
\[
                   P(v)=-3\bigl(9-3(1+1+1)\bigr)=0.
\]
For every $r\in\mathbb Z$ and $\varepsilon\in\{-1,1\}$, put
\[
 L_\varepsilon(r)=(-3r-\varepsilon,r,r,r+\varepsilon)
                =rv+(-\varepsilon,0,0,\varepsilon).
 \label{d4i:eq:two-lines}\tag{D4I7}
\]
Its coordinate sum is
$-3r-\varepsilon+r+r+r+\varepsilon=0$, so it belongs to $D_4$.
The original last-three-coordinate square sum and first-coordinate
square are respectively
\[
 s=r^2+r^2+(r+\varepsilon)^2=3r^2+2r\varepsilon+1,
 \qquad (-3r-\varepsilon)^2=9r^2+6r\varepsilon+1.
\]
Thus the factor in the original cubic is exactly
\[
 a_0^2-3s=(9r^2+6r\varepsilon+1)
                 -(9r^2+6r\varepsilon+3)=-2,
\]
and therefore
\[
 P(L_\varepsilon(r))=(-3r-\varepsilon)(-2)=6r+2\varepsilon.
 \label{d4i:eq:two-line-values}\tag{D4I8}
\]

For every $h\in\mathbb Z$, define the third line
\[
 L_0(h)=(-3h,h+1,h-1,h)=hv+(0,1,-1,0).
 \label{d4i:eq:third-line}\tag{D4I9}
\]
Its coordinate sum is zero, so it also belongs to $D_4$. Here
\[
 s=(h+1)^2+(h-1)^2+h^2=3h^2+2,
\]
and consequently
\[
 P(L_0(h))=(-3h)\bigl(9h^2-3(3h^2+2)\bigr)
                         =(-3h)(-6)=18h.
 \label{d4i:eq:third-line-values}\tag{D4I10}
\]
These identities hold for all signed integer parameters. In particular,
$L_0(0)=(0,1,-1,0)$ has cubic value zero, and the zero vector also
has cubic value zero.

\subsection{The global image and a proved reconstruction map}
\label{d4i:sec:global-image}

Set
\[
 S=\{n\in\mathbb Z:n\in2\mathbb Z\text{ and }
                                  (3\nmid n\text{ or }9\mid n)\}.
 \label{d4i:eq:image-domain}\tag{D4I11}
\]

\begin{theorem}[Complete cubic image and explicit section]
\label{d4i:thm:complete-image}
The original polynomial has the exact image
\[
 P(D_4)=18\mathbb Z\ \cup\ (6\mathbb Z+2)\ \cup\ (6\mathbb Z-2)=S.
 \label{d4i:eq:complete-image}\tag{D4I12}
\]
Equivalently,
\[
 P(D_4)=\{n\in\mathbb Z:
                 n\bmod18\in\{0,2,4,8,10,14,16\}\}.
 \label{d4i:eq:residue-image}\tag{D4I13}
\]
The following map is an actual section $\sigma:S\to D_4$:
\[
 \sigma(n)=
 \begin{cases}
 (-3h,h+1,h-1,h),&n=18h,\\
 (-3r-\varepsilon,r,r,r+\varepsilon),
        &n=2(3r+\varepsilon),\quad\varepsilon\in\{-1,1\}.
 \end{cases}
 \label{d4i:eq:global-section}\tag{D4I14}
\]
Its displayed parameters are uniquely determined in their respective
cases, and $P\circ\sigma=\operatorname{id}_S$.
\end{theorem}

\begin{proof}
Equation~\eqref{d4i:eq:necessary-restrictions} proves
$P(D_4)\subseteq S$. For the reverse inclusion let $n\in S$.
If $3\mid n$, then $9\mid n$, and evenness gives $n=18h$ for
a unique $h\in\mathbb Z$. Equation~\eqref{d4i:eq:third-line-values}
gives $P(L_0(h))=n$.

If $3\nmid n$, write $n=2q$. Then $3\nmid q$, so there is a
unique $\varepsilon\in\{-1,1\}$ with
$q\equiv\varepsilon\pmod3$. The number
$r=(q-\varepsilon)/3$ is consequently a unique integer.
Equation~\eqref{d4i:eq:two-line-values} gives
\[
                    P(L_\varepsilon(r))=6r+2\varepsilon=2q=n.
\]
The two displayed cases are disjoint and cover $S$, and both
witnesses lie in $D_4$ by \cref{d4i:sec:three-lines}. These
calculations prove the asserted section identity and both inclusions
in \eqref{d4i:eq:complete-image}. The first arithmetic progression
contributes residue zero modulo eighteen; the second contributes
$2,8,14$; and the third contributes $16,4,10$. This proves
\eqref{d4i:eq:residue-image}.
\end{proof}

The classification includes every positive, zero, and negative value.
The full excluded set is the union of the odd integers with
$18\mathbb Z+6$ and $18\mathbb Z+12$.
Let
\[
 A=\{a\in\mathbb Z^4:a_0+a_1+a_2+a_3=0\}
 \label{d4i:eq:sum-zero-slice}\tag{D4I15}
\]
be the literal sum-zero sublattice in the original coordinates. The
three witnesses lie in $A\subseteq D_4$, so the two inclusions
$S\subseteq P(A)\subseteq P(D_4)=S$ prove
\[
                              P(A)=P(D_4)=S.
 \label{d4i:eq:slice-image}\tag{D4I16}
\]
The image of the original cubic $F=\tau$ on the common lattice is
therefore exactly
\[
\begin{aligned}
 F(\mathcal I)&=2\sqrt2\,S\\
 &=4\sqrt2\bigl((9\mathbb Z)\cup(3\mathbb Z+1)
                                     \cup(3\mathbb Z-1)\bigr).
\end{aligned}
 \label{d4i:eq:original-cubic-image}\tag{D4I17}
\]
Because the original off-diagonal Albert determinant is $2F$, its
corresponding image is $4\sqrt2\,S$. The additive group generated
by that determinant image is $8\sqrt2\mathbb Z$, while the
additive group generated by $F(\mathcal I)$ is $4\sqrt2\mathbb Z$.
Indeed every $P$ value is even, and the first two lines at parameter
zero attain $P=2$ and $P=-2$. Thus the generated additive group of
the $P$ values is exactly $2\mathbb Z$; applying the two original
scalar factors proves both group assertions without identifying an
image set with its generated group.

\subsection{Exact local images with explicit local sections}
\label{d4i:sec:local-images}

For a prime $p$, let
\[
 D_{4,p}=\{a\in\mathbb Z_p^4:
                            a_0+a_1+a_2+a_3\in2\mathbb Z_p\}.
 \label{d4i:eq:local-domain}\tag{D4I18}
\]
This is precisely $D_4\otimes_{\mathbb Z}\mathbb Z_p$ in the original
coordinates. To verify the equality directly, use the integer basis
\[
\begin{gathered}
 b_0=(1,-1,0,0),\quad b_1=(0,1,-1,0),\\
 b_2=(0,0,1,-1),\quad b_3=(0,0,1,1).
\end{gathered}
 \label{d4i:eq:integral-basis}\tag{D4I19}
\]
For any $a$ in the displayed definition of $D_{4,p}$, the unique
coefficients of these basis vectors are
\[
\begin{aligned}
 c_0&=a_0,&c_1&=a_0+a_1,\\
 c_3&=(a_0+a_1+a_2+a_3)/2,&c_2&=c_3-a_3.
\end{aligned}
 \label{d4i:eq:local-basis-inverse}\tag{D4I20}
\]
All belong to $\mathbb Z_p$, and substitution gives the original
four coordinates. Conversely a $\mathbb Z_p$-linear combination
of the four vectors has coordinate sum $2c_3$. These two direct
calculations prove the tensor-product identification, including
at $p=2$. Taking integral coefficients gives the same basis
description of the original $D_4$.

The three line identities of \cref{d4i:sec:three-lines} hold over
$\mathbb Z_p$, since they are integer polynomial identities. Each
line stays in the sum-zero submodule of $D_{4,p}$.

\begin{theorem}[All local images and the exact local-to-global statement]
\label{d4i:thm:local-images}
For the unchanged polynomial $P$,
\[
\begin{aligned}
 P(D_{4,2})&=2\mathbb Z_2,\\
 P(D_{4,3})&=\mathbb Z_3^{\times}\cup9\mathbb Z_3,\\
 P(D_{4,p})&=\mathbb Z_p\qquad(p\ge5).
\end{aligned}
 \label{d4i:eq:local-images}\tag{D4I21}
\]
An integer $n$ belongs to every local image if and only if
$n\in S$, and hence if and only if $n\in P(D_4)$. The global
section is exactly \eqref{d4i:eq:global-section}; explicit local
sections are given in the proof.
\end{theorem}

\begin{proof}
For $p=2$, the parity calculation of \cref{d4i:sec:necessary}, reduced
modulo two, gives $P(D_{4,2})\subseteq2\mathbb Z_2$. Conversely,
given an arbitrary $n=2q\in2\mathbb Z_2$, put
\[
 r=(q-1)/3\in\mathbb Z_2.
\]
Division by three is valid since three is a unit in $\mathbb Z_2$.
Then $L_1(r)\in D_{4,2}$ and
\[
                         P(L_1(r))=6r+2=2q=n.
\]
Thus $n\mapsto L_1((n/2-1)/3)$ is an explicit section on
$2\mathbb Z_2$, and proves the first equality.

For $p=3$, reduction of the original polynomial modulo three and
the substitution $a_0=3b$ in
\eqref{d4i:eq:three-divisibility} prove that any value divisible by
three belongs to $9\mathbb Z_3$. Every unit
$n\in\mathbb Z_3^{\times}$ is congruent modulo three to exactly
one of $2\varepsilon$, with $\varepsilon\in\{-1,1\}$. Set
\[
 r=(n-2\varepsilon)/6\in\mathbb Z_3.
\]
This quotient is integral because the numerator is divisible by
three and two is a unit. The line $L_\varepsilon(r)$ then has
$P$ value $6r+2\varepsilon=n$.
Every $n\in9\mathbb Z_3$ can be written uniquely as $18h$ with
$h=n/18\in\mathbb Z_3$, again because two is a unit.
The third line has $P(L_0(h))=18h=n$. These two disjoint cases
give the explicit local section
\[
 n\longmapsto
 \begin{cases}
 L_\varepsilon((n-2\varepsilon)/6),
        &n\in\mathbb Z_3^{\times},\quad n\equiv2\varepsilon\pmod3,\\
 L_0(n/18),&n\in9\mathbb Z_3.
 \end{cases}
 \label{d4i:eq:three-adic-section}\tag{D4I22}
\]
They prove the second equality in \eqref{d4i:eq:local-images}.

For every prime $p\ge5$, six is a unit. Given any
$n\in\mathbb Z_p$, choose $r=(n-2)/6\in\mathbb Z_p$.
The point $L_1(r)$ belongs to $D_{4,p}$ and has
$P(L_1(r))=6r+2=n$. Thus $n\mapsto L_1((n-2)/6)$ is a
section on all of $\mathbb Z_p$, proving the third equality.

For integers $n$, the first local condition is precisely evenness;
the second is precisely $3\nmid n$ or $9\mid n$; and the other
primes impose no condition. Membership in every local image is
therefore exactly $n\in S$. The explicit integral section in
\cref{d4i:thm:complete-image} proves that those conditions suffice
for global integral representability. This proves the local-to-global
assertion without approximation or inference from a finite modular
calculation.
\end{proof}

\subsection{The exact factorization and quotient metric on the sum-zero slice}
\label{d4i:sec:factorization}

Define the integral map into the original sublattice $A$ of
\eqref{d4i:eq:sum-zero-slice} by
\[
 \Phi:\mathbb Z^3\longrightarrow A,\qquad
 \Phi(w,u,v)=(-3w-u-v,w+u,w+v,w).
 \label{d4i:eq:slice-map}\tag{D4I23}
\]
Its coordinate sum is zero. For $a\in A$, its proposed inverse is
\[
                 \Phi^{-1}(a)=(a_3,a_1-a_3,a_2-a_3).
 \label{d4i:eq:slice-inverse}\tag{D4I24}
\]
Substitution returns the original second, third, and fourth
coordinates immediately; the first becomes $-a_1-a_2-a_3=a_0$
by the defining equation of $A$. Substituting $\Phi$ into the
inverse returns $(w,u,v)$. These two equalities prove that $\Phi$
is an exact integral-module isomorphism.

For $a=(a_0,y,z,w)\in A$, put $T=y+z+w=-a_0$. Direct expansion
of the original polynomial gives
\[
\begin{aligned}
 P(a)&=-T\bigl(T^2-3(y^2+z^2+w^2)\bigr)\\
     &=2T(y^2+z^2+w^2-yz-yw-zw)\\
     &=T\bigl((y-z)^2+(y-w)^2+(z-w)^2\bigr).
\end{aligned}
 \label{d4i:eq:slice-expansion}\tag{D4I25}
\]
Under $\Phi$, one has $T=3w+u+v$, and the three differences are
$u-v,u,v$. Consequently
\[
                 P(\Phi(w,u,v))=2(3w+u+v)(u^2-uv+v^2).
 \label{d4i:eq:factorization}\tag{D4I26}
\]
This identity explains the three original lines without changing
their coordinates. On $L_\varepsilon(r)$, the quotient differences
are $(u,v)=(-\varepsilon,-\varepsilon)$, so
$u^2-uv+v^2=1$; on $L_0(h)$, they are $(u,v)=(1,-1)$, so
$u^2-uv+v^2=3$. The original first coordinates remain
$-3r-\varepsilon$ and $-3h$, respectively.

The exact quotient map is
\[
 \chi:A\longrightarrow\mathbb Z^2,\qquad
                  \chi(a)=(a_1-a_3,a_2-a_3).
 \label{d4i:eq:quotient-map}\tag{D4I27}
\]
It is onto because $\chi\Phi(0,u,v)=(u,v)$. An element lies in
its kernel exactly when its last three coordinates all equal a
common integer $w$; its first coordinate is then $-3w$.
Thus $\ker\chi=\mathbb Z(-3,1,1,1)$. The section
$(u,v)\mapsto\Phi(0,u,v)$ proves the split exact sequence
\[
 0\longrightarrow\mathbb Z(-3,1,1,1)\longrightarrow A
       \xrightarrow{\ \chi\ }\mathbb Z^2\longrightarrow0.
 \label{d4i:eq:split-exact-sequence}\tag{D4I28}
\]

For $\omega=(-1+\sqrt{-3})/2$, one has
$\omega+\overline\omega=-1$ and $\omega\overline\omega=1$.
Direct multiplication therefore proves
\[
                 (u+v\omega)(u+v\overline\omega)=u^2-uv+v^2.
 \label{d4i:eq:eisenstein-norm}\tag{D4I29}
\]
The quadratic factor in the original cubic is exactly this
Eisenstein norm, with the preceding coefficient two unchanged.

The original Euclidean metric on these same coordinates is,
by expansion,
\[
 \|\Phi(w,u,v)\|^2
            =12w^2+8w(u+v)+2(u^2+uv+v^2).
 \label{d4i:eq:full-slice-metric}\tag{D4I30}
\]
Write $d=(-3,1,1,1)$ for the direction when used in the metric
calculation. Its squared norm is twelve, and the exact pairing is
\[
          \langle\Phi(w,u,v),d\rangle=12w+4(u+v).
 \label{d4i:eq:direction-pairing}\tag{D4I31}
\]
Hence orthogonal projection away from the real line $\mathbb R d$
is given in the original coordinates by
\[
\begin{aligned}
 \pi\Phi(w,u,v)
  &=\Phi(w,u,v)-\left(w+\frac{u+v}{3}\right)d\\
  &=\left(0,\frac{2u-v}{3},\frac{2v-u}{3},-\frac{u+v}{3}\right).
\end{aligned}
 \label{d4i:eq:orthogonal-projection}\tag{D4I32}
\]
This depends only on the quotient coordinate $\chi(a)=(u,v)$.
Its squared norm is
\[
\begin{aligned}
 \|\pi\Phi(w,u,v)\|^2
  &=\frac{(2u-v)^2+(2v-u)^2+(u+v)^2}{9}\\
  &=\frac23(u^2-uv+v^2).
\end{aligned}
 \label{d4i:eq:projected-norm}\tag{D4I33}
\]
The quotient metric on $A/\mathbb Z d$, realized by this exact
projection, thus has Gram matrix
\[
              \frac13\begin{pmatrix}2&-1\\-1&2\end{pmatrix}
 \label{d4i:eq:quotient-gram}\tag{D4I34}
\]
in the displayed quotient coordinates $(u,v)$.
The retained common-lattice map $a\mapsto\sqrt2(a,a,a)$ multiplies
every squared norm by six, as proved in
\cref{wcl:thm:literal-intersection}. Consequently the corresponding
quotient Gram matrix on the image of $A$ in $\mathcal I$ is
\[
                         \begin{pmatrix}4&-2\\-2&4\end{pmatrix}.
 \label{d4i:eq:common-lattice-quotient-gram}\tag{D4I35}
\]
This is the exact quotient-metric calculation on the indicated
rank-three slice. It neither identifies that slice with the full
$D_4$ nor changes the original common-lattice metric.
